% Einheit 2 von 4 – Satz vom Nullprodukt (bank/quadratische-gleichungen/e2.jsonl, zusammenbau v0.1)
%% TODO Zweigzeile Teil 1: was hier gelernt wird (Satz in Schülersprache aus dem Einheitstitel) – Einheit 2
\einheitenkopf[e2]{Einheit 2 von 4 · Satz vom Nullprodukt}
\zweigzeile{ab Kl. 10 (am Gymnasium ab Kl. 8) · P10}
\setcounter{aufgabe}{13}

%% TODO Titel: Ich-kann-Satz fehlt in der Bank (Platzhalter: Kettenname) – Nr. 14
%% TODO Anweisung: ein Satz über der ersten Teilaufgabe fehlt in der Bank – Nr. 14
\begin{aufgabe}{Nullprodukt}
%% TODO \rechenplatz{n}: Schrittzahl des Grundfalls fehlt in der Bank (nur bei fester Schreibform, 2.3 b)
\begin{teile}
\teil $(x - 2) \cdot (x + 7) = 0$ – gilt der Satz vom Nullprodukt?\\ \kreuz{gilt} \\ \kreuz{gilt nicht}
\end{teile}
\begin{gleichungsraster}[2]
\gl{\text{(x - 2) \cdot (x - 5) = 0}} &
\gl{\text{(x + 4) \cdot (x - 1) = 0}} \\
\gl{\text{x \cdot (x + 6) = 0}} &
\gl{\text{(x - 7) \cdot (x - 7) = 0}} \\
\end{gleichungsraster}
\begin{teile}
\teil Ist $x = -5$ Lösung von $(x + 5) \cdot (x - 1) = 0$? \janein
\end{teile}
\begin{gleichungsraster}[2]
\gl{\text{x^2 + 9x = 0}} &
\gl{\text{x^2 - 3x = 0}} \\
\gl{\text{2x^2 - 6x = 0}} &
\end{gleichungsraster}
\begin{teile}
\teil $x \cdot (x + 3) = 10$ – ausmultiplizieren und ordnen? \leerfeld = 0
\teil Welcher Wert erfüllt $x \cdot (x + 9) = -14$? \hfill (P10 2025 OS)\\ \kreuz{$x = 2$} \\ \kreuz{$x = 7$} \\ \kreuz{$x = -7$} \\ \kreuz{$x = -14$}
\end{teile}
\end{aufgabe}

%% TODO Titel: Ich-kann-Satz fehlt in der Bank (Platzhalter: Kettenname) – Nr. 15
%% TODO Anweisung: ein Satz über der ersten Teilaufgabe fehlt in der Bank – Nr. 15
\begin{aufgabe}{Nullprodukt – Fehler finden, Begründen}
\begin{teile}
\teil Wo ist der Fehler? \rechnung{x^2 &= 6x &&\mid : x \\ x &= 6}
\teil Begründe: Ist bei $(x - 3) \cdot (x + 5)$ einer der beiden Faktoren null, ist das ganze Produkt null.
\end{teile}
\end{aufgabe}
